\section{Failure Case Analysis}
\label{sec:failure}

\Cref{fig:fail} shows two scenarios where Canny --- a pure first-order-gradient
detector with no notion of semantics --- fails. In both, crack recall stays at
\num{1.0} but precision collapses ($P=\num{0.41}$ and $P=\num{0.35}$).

\paragraph{Strong cast-shadow contour}
A hard shadow boundary is a genuine intensity step, so it produces a Sobel
response \emph{larger} than the crack; NMS keeps it and no $(\Thi,\Tlo)$ removes
it without also removing the crack. The limitation is definitional: Canny
detects \emph{intensity} edges, but a shadow is an \emph{illumination} edge, not
surface geometry. Separating them needs illumination modelling (retinex /
homomorphic) or colour, which a single-channel gradient lacks.

\paragraph{High-frequency texture}
On coarse aggregate the surface is itself full of small intensity steps whose
magnitudes match the faint crack, so (i) no global threshold separates signal
from texture, and (ii) hysteresis worsens it --- the texture edges form a
percolating 8-connected web, so one strong seed floods large areas. The output
is an "edge soup" (73 components) in which the crack is not distinguishable by
connectivity. The limitation is scale: a fixed-$\sigma$ gradient cannot tell
which edges belong to a structure of interest. Remedies: oriented crack-matched
filters, multi-scale ridge measures (Frangi), or a learned segmenter.

\begin{figure}[tb]
  \centering
  \includegraphics[width=0.7\linewidth]{task5_failure_cases.png}
  \caption{Failure cases. Row 1: a cast-shadow contour becomes a strong false
  edge. Row 2: dense texture yields a fragmented map that buries the crack.
  Columns: input, Canny output, ground truth.}
  \label{fig:fail}
\end{figure}
